% !TEX root = ../main.tex
\section{Related Work}

\subsection{UAV-Based Wildfire Rescue}

UAVs provide a practical means of observing wildfires in hazardous and
difficult-to-access areas, but operational wildfire rescue requires more than
detecting fire in aerial imagery. It requires a decision chain that transforms
observations into a current representation of the incident, allocates limited
airborne resources to the resulting tasks, and updates the mission when the
environment or fleet state changes. Two properties distinguish this problem
from a collection of independent UAV capabilities. Decisions at one stage
constrain the next stage, and decisions that are valid initially may cease to be
valid during execution. Existing work has made substantial progress on the
individual stages of this chain, including aerial perception, active situation
monitoring, autonomous flight planning, and multi-UAV coordination, but the
connection among these stages remains less developed.

At the perception stage, multimodal sensing improves the reliability of fire
observations under the visually challenging conditions common in wildfire
scenes. Thermal--RGB fusion has been used to detect fire and hotspots from UAV
imagery, reducing sensitivity to smoke and illumination variation
\cite{yuan2021uavfire}. Transformer-based aerial semantic segmentation further
supports the interpretation of burned regions, vegetation, and terrain
\cite{zhang2022uavtransformer}, while learning-based 3D reconstruction can
provide geometric environmental models for navigation and inspection
\cite{yang2022uav3dreconstruction}. These methods produce the visual and
spatial evidence needed for rescue, but their outputs must still be converted
into mission-level entities, such as fire points, task zones, and safety-relevant
context, before they can drive rescue actions. This conversion is consequential:
perception determines what the planner treats as a task, yet perception quality
alone does not indicate whether the resulting tasks can be assigned to a
resource-limited fleet or sustained throughout execution.

Situation awareness and planning research addresses part of this conversion.
Autonomous exploration and mapping enable UAVs to update knowledge of unknown
or changing disaster environments \cite{popovic2021uavdisaster}; next-best-view
planning selects viewpoints that maximize information gain and reduces
observation uncertainty \cite{bircher2016nbv}. On the action side, deep
reinforcement learning has been applied to autonomous navigation in unknown
environments \cite{xiao2021uavrl}, and coverage path planning supports
large-area search and perimeter scanning \cite{yu2022coverageuav}. However,
these methods generally optimize sensing, navigation, or coverage as an
individual capability. They do not by themselves produce a fleet-level,
resource-constrained mission plan that explicitly links evolving fire
information to task priorities, payload use, endurance, and execution status.
In particular, optimizing where a UAV should observe or fly is different from
deciding which combination of vehicles, routes, and payloads should serve
multiple interdependent zones under a common mission objective.

Multi-UAV methods extend decision making from a single vehicle to a team.
Multi-agent reinforcement learning has demonstrated cooperative navigation and
task allocation capabilities \cite{yang2023multiagentuav}, while the broader
disaster-response literature highlights the value of UAV support and
human--robot collaboration in search-and-rescue operations
\cite{waharte2021uavsar}. Nevertheless, wildfire rescue remains difficult
because fire evolution, communication changes, energy depletion, and vehicle
failures can invalidate an initially feasible allocation. The literature is
therefore still dominated by modular or open-loop solutions: perception,
mapping, path planning, and coordination are often evaluated separately, and
task allocation is commonly assessed from a fixed initial state. This leaves a
gap between initial coordination and mission continuity: a method may produce a
valid allocation at one instant without specifying how the fleet should recover
when the task set, vehicle state, or resource constraints change.

The methods reviewed above also assume different inputs, outputs, resource
models, and evaluation protocols. Their reported results provide useful
references for individual capabilities, but do not constitute a directly
comparable end-to-end numerical baseline for a task that begins with aerial
observations and continues through resource-constrained mission planning and
runtime replanning.
The relevant gap is therefore not the absence of perception, planning, or
coordination algorithms in isolation. It is the absence of a common high-level
decision loop that connects these capabilities and exposes how resource
constraints and dynamic events propagate from fire-situation understanding to
fleet execution.

\subsection{LLM-Based Robotic Agents}

Large language models (LLMs) provide an increasingly capable interface between
high-level human intent and robot behavior. Rather than replacing low-level
perception or control modules, they are commonly used to interpret language,
decompose goals, select skills, and reason over structured observations. This
role is particularly relevant to wildfire rescue, where a decision system
must translate a changing fire situation and operational constraints into a
coherent fleet mission rather than execute a fixed motion policy. For such a
role, generative flexibility is only one requirement. The resulting decisions
must also be grounded in available capabilities, checked against operational
constraints, and revised when feedback invalidates their assumptions.

Language grounding and hierarchical planning are central directions in this
area. SayCan combines language-model scoring with robot affordance estimates so
that high-level instructions are grounded in executable skills
\cite{ahn2022saycan}. Code as Policies uses language models to generate
programmatic control policies, providing a flexible mechanism for composing
robot behavior from functions and APIs \cite{liang2023codeaspolicies}.
Related design guidance has shown how ChatGPT can be used to generate robot
workflows while exposing tool descriptions and environmental constraints to the
model \cite{vemprala2023chatgptrobotics}. Together, these approaches show that
LLMs can support task decomposition and action sequencing, but they also make
clear that generated actions require grounding and external feasibility checks.
This distinction is central in multi-UAV missions: linguistic coherence does
not guarantee that a set of assignments respects endurance, payload,
communication, priority, and safety constraints simultaneously.

Embodied and multimodal models extend this reasoning interface to visual and
physical state. RT-1 learns vision-language-action policies from large-scale
real-robot data \cite{brohan2023rt1}, while PaLM-E incorporates visual inputs
and robot states into an embodied language model \cite{driess2023palme}.
Socratic Models further demonstrate zero-shot reasoning by composing
specialized vision and language models \cite{zeng2023socratic}. These methods
strengthen the connection between perception and action, yet they are generally
evaluated on individual robots and predefined manipulation or navigation tasks.
They do not directly address resource-constrained mission planning for dynamic,
interdependent tasks across a rescue fleet. Moving from an individual embodied
agent to a fleet-level decision layer introduces coupling among vehicle states,
task priorities, routes, and shared resources, so an action that is feasible for
one vehicle may still be unsuitable for the mission as a whole.

LLM-based agents have also explored closed-loop decision making through
environmental feedback, navigation, and accumulated experience. Inner Monologue
uses language-based feedback to revise embodied plans
\cite{huang2023innermonologue},
and LM-Nav combines large pretrained models with semantic waypoints for
open-world navigation \cite{shah2023lmnav}. Voyager couples an LLM agent with
an automatic curriculum and a growing skill library for open-ended exploration
\cite{wang2023voyager}. Nevertheless, these systems primarily target single-
agent embodied environments. They establish that feedback and accumulated
experience can improve sequential behavior, but do not establish how these
ideas should be combined with fleet-wide resource constraints or how their
importance changes as operational pressure increases in multi-UAV missions.

Taken together, the literature provides the main ingredients for high-level
robot decision making, but not their integration around the specific failure
modes of dynamic, resource-constrained fleet missions. An LLM-centered
closed-loop decision framework for this setting must satisfy three requirements
beyond generating an initial plan. It must reject plans that violate operational
or tactical constraints, retrieve relevant experience when resource and scheduling
interactions make planning less routine, and repair the active mission when
execution feedback changes the underlying state. These requirements motivate a
closed-loop framework in which Validation, Memory, and Runtime Replanning are
not independent additions, but complementary mechanisms for sustaining
high-level decision making from initial observation through execution.

\endinput
